\documentclass[11pt]{article}
\usepackage[margin=1in]{geometry}
\usepackage{listings}
\usepackage[hidelinks]{hyperref}
\usepackage{enumitem}
\lstdefinelanguage{JavaScript}{
  keywords={const, let, var, function, return, if, else, for, while, class, new, this, extends, static},
  sensitive=true,
  morecomment=[l]{//},
  morecomment=[s]{/*}{*/},
  morestring=[b]",
  morestring=[b]'
}
\lstset{
  basicstyle=\ttfamily\small,
  columns=fullflexible,
  frame=single,
  breaklines=true
}

\title{Gruggle JavaScript API Guide}
\author{}
\date{\today}

\begin{document}
\maketitle

\section{Overview}
Gruggle is a single-file DOT loader/transformer that can be used both as a CLI and as a programmatic API from Node.js. The CLI remains unchanged; when \texttt{gruggle.js} is imported via \texttt{require}, no CLI code runs and the API is exposed.

\section{Getting Started}
\begin{itemize}[leftmargin=1.5em]
  \item Ensure Node.js is available.
  \item Place \texttt{gruggle.js} in your project (or install it alongside your code).
  \item Use either the CLI (\texttt{node gruggle.js ...}) or import it:
\end{itemize}
\begin{lstlisting}[language=JavaScript]
const {
  Gruggle,
  parse,
  serialize,
  applySubcommands,
  applyOne,
  parseChainString,
  emptyGraph,
  mergeGraphs,
} = require('./gruggle');
\end{lstlisting}

\section{CLI vs Library}
\begin{itemize}[leftmargin=1.5em]
  \item CLI: invoked via \texttt{node gruggle.js ...}; supports all documented subcommands.
  \item Library: when required, the module exports API objects; \texttt{main()} is not executed (checked via \texttt{require.main === module}).
\end{itemize}

\section{Core Concepts}
\begin{itemize}[leftmargin=1.5em]
  \item \textbf{Graph object}: plain JS object containing nodes (Map), edges (array), defaults, and attributes.
  \item \textbf{Subcommands}: the same operations used by the CLI; they can be invoked programmatically via \texttt{applyOne} or \texttt{applySubcommands}.
  \item \textbf{Chain strings}: CLI-style quoted command strings parsed by \texttt{parseChainString}.
\end{itemize}

\section{Class: \texttt{Gruggle}}
Convenience wrapper around the core helpers.

\subsection*{Construction}
\begin{description}[leftmargin=1.8em, style=nextline]
  \item[\texttt{new Gruggle(graph = emptyGraph())}] Start from an existing graph object or a fresh empty graph.
  \item[\texttt{Gruggle.fromDot(text, edgeMode = 'curved')}] Parse DOT text; \texttt{edgeMode} can be \texttt{'curved'} (default) or \texttt{'ortho'}/\texttt{'rectilinear'} (sets \texttt{splines=ortho}).
  \item[\texttt{Gruggle.fromFiles(files, edgeMode = 'curved')}] Parse and merge a list of DOT files, then set edge mode.
\end{description}

\subsection*{Methods}
\begin{description}[leftmargin=1.8em, style=nextline]
  \item[\texttt{chain(commands)}] Run one or more subcommands using the \texttt{chain} wrapper. Accepts a single chain string (e.g., \texttt{"add-nodes n:1:3"}) or an array of chain strings or token arrays. Returns \texttt{this}.
  \item[\texttt{apply(tokens)}] Apply a token list directly (e.g., \texttt{['add-nodes', 'n:1:3']}); returns \texttt{this}.
  \item[\texttt{toDot()}] Serialize the current graph to DOT text (obeys keep/remove flags).
  \item[\texttt{toSvg(outFile)}] Emit SVG via Graphviz \texttt{dot}, writing to \texttt{outFile}; returns \texttt{this}.
  \item[\texttt{value()}] Return the underlying graph object (for advanced inspection/manipulation).
\end{description}

\section{Helper Functions}
\begin{description}[leftmargin=1.8em, style=nextline]
  \item[\texttt{parse(text)}] Parse DOT text into a graph object.
  \item[\texttt{serialize(graph)}] Serialize a graph to DOT text, honoring keep/remove semantics.
  \item[\texttt{applyOne(graph, tokens)}] Apply a single subcommand token array to a graph.
  \item[\texttt{applySubcommands(graph, tokens)}] Apply a token array (often starting with a subcommand, e.g., \texttt{['add-nodes', 'n:1:3']} or \texttt{['chain', 'cmd1', 'cmd2']}).
  \item[\texttt{parseChainString(str)}] Parse a CLI-style chain string into tokens, handling quoted segments.
  \item[\texttt{emptyGraph()}] Create a new empty digraph with defaults.
  \item[\texttt{mergeGraphs(graphs)}] Merge multiple graphs (as the CLI does for multiple \texttt{-i} inputs).
\end{description}

\section{Usage Examples}
\subsection*{Build and serialize}
\begin{lstlisting}[language=JavaScript]
const { Gruggle } = require('./gruggle');
const gm = new Gruggle();
gm.apply(['add-nodes', 'n:1:3'])
  .apply(['add-edges', 'lbl', 'n.*', 'n.*']);
console.log(gm.toDot());
\end{lstlisting}

\subsection*{Chain with CLI-style strings}
\begin{lstlisting}[language=JavaScript]
const { Gruggle, parseChainString } = require('./gruggle');
const gm = Gruggle.fromFiles(['graph1.g']);
gm.chain(parseChainString('add-nodes extra:1:2'))
  .chain(['add-edges', 'note', 'extra.*', 'n.*'])
  .toSvg('out.svg');
\end{lstlisting}

\subsection*{Direct helpers without Gruggle}
\begin{lstlisting}[language=JavaScript]
const { parse, applySubcommands, serialize, parseChainString } = require('./gruggle');
let g = parse('digraph G { a -> b; }');
g = applySubcommands(g, ['add-nodes', 'c']);
g = applySubcommands(g, ['add-edges', 'lbl', 'a.*', 'c']);
console.log(serialize(g));
\end{lstlisting}

\section{Notes and Compatibility}
\begin{itemize}[leftmargin=1.5em]
  \item The CLI behavior is unchanged; all existing tests and subcommands remain valid.
  \item The API lives in the same single file; distribution remains a single \texttt{gruggle.js}.
  \item \texttt{toSvg} requires Graphviz \texttt{dot} on PATH, same as the CLI.
  \item Chain parsing rules match the CLI; use \texttt{parseChainString} when you want to feed a quoted command string.
\end{itemize}

\end{document}
